% !TEX root = ../main.tex
% This file is \input{}-ed by main.tex, so it must NOT contain \documentclass, package imports,
% or \begin{document}/\end{document}. Put summary content only.

\begin{center}
    {\LARGE\bfseries Forecasts for Smartphone Battery Discharge Dynamics\\[0.5ex] and User Behavior Optimization\par}
    \vspace{1.5ex}
    {\bfseries Summary\par}
\end{center}
\vspace{1.5ex}

As the primary power source for modern smartphones, lithium-ion batteries directly determine device usability and user experience. Due to complex electrochemical processes and varying usage patterns, accurately predicting battery life remains a challenging task. The aim of this report is to build a comprehensive battery discharge prediction model to evaluate the influences of various factors on battery life. We are expected to provide actionable strategies for both users and operating system designers to extend battery life under the constraints of practical usage conditions. Three models are established: \textbf{Model I: SOC Continuous-Time Dynamics Model}; \textbf{Model II: Time-to-Empty Prediction Model}; \textbf{Model III: Performance Impact and Aging Model}.

\textbf{For Model I}, we establish a five-state continuous-time ordinary differential equation (ODE) system based on a second-order Equivalent Circuit Model (2-RC ECM). The state vector includes State of Charge (SOC), two RC-branch voltages ($V_{RC1}$, $V_{RC2}$), battery temperature ($T$), and State of Health (SOH). The Open Circuit Voltage (OCV) is modeled using a Shepherd-enhanced Nernst equation derived from electrochemical principles. The core SOC dynamics follow Coulomb counting:
\begin{equation*}
    \frac{d\SOC}{dt} = -\frac{I(t)}{\Qnom \cdot \SOH \cdot 3600}
\end{equation*}
The current $I(t)$ is computed from power consumption through KVL analysis of the equivalent circuit. The results demonstrate that this electrochemistry-based approach achieves RMSE of 2.1\% for SOC estimation across standard drive cycles, as shown in Table~2.

\textbf{For Model II}, building upon Model I, we develop an energy-based time-to-empty prediction framework. The basic analytical formula is:
\begin{equation*}
    t_{\text{empty}} = \frac{Q_{\text{eff}}(\SOH, T) \cdot (\SOC_0 - \SOC_{\min}) \cdot V_{\text{avg}}}{P_{\text{avg}}}
\end{equation*}
Temperature correction factors are introduced to account for capacity reduction at extreme temperatures. Monte Carlo simulation with 10,000 samples is employed to quantify prediction uncertainty, yielding 95\% confidence intervals. The simulation results, presented in Figures~1--3, show that battery life varies from 2.1 hours (heavy gaming at low temperature) to 22.5 hours (light usage at optimal temperature), representing a 10.7$\times$ dynamic range.

\textbf{For Model III}, we model battery aging through two mechanisms: calendar aging and cycle aging. Calendar aging follows an Arrhenius-type temperature dependence with SOC stress factor:
\begin{equation*}
    \frac{d\SOH_{\text{cal}}}{dt} = -A_{\text{cal}} \cdot \exp\left(-\frac{E_a}{RT}\right) \cdot \exp\bigl(\beta \cdot (\SOC - 0.5)\bigr)
\end{equation*}
Cycle aging depends on depth of discharge (DOD) and C-rate:
\begin{equation*}
    \frac{d\SOH_{\text{cyc}}}{dt} = -k_{\text{cyc}} \cdot \text{DOD}^{\gamma} \cdot \left(\frac{I}{I_{\text{nom}}}\right)^{\alpha} \cdot |I|
\end{equation*}
The combined aging model predicts that under ``Gaming + Fast Charge'' conditions (high temperature, high SOC, high current), battery capacity degrades to 80\% (End-of-Life threshold) within approximately 1.3 years, while ``Light Use'' scenarios extend battery life to over 3 years before reaching EOL, as illustrated in Figure~9.

\textbf{In addition}, this report conducts Sobol variance-based global sensitivity analysis to identify the key parameters influencing battery life predictions. Based on 2,560 model evaluations with Saltelli's sampling scheme, the results in Table~3 reveal that: Initial SOC dominates the output variance ($S_T = 0.71$); CPU utilization is the primary controllable factor ($S_T = 0.39$); Screen brightness has moderate impact ($S_T = 0.08$); Temperature effects are surprisingly negligible within normal operating range ($S_T < 0.01$).

\textbf{Regarding user recommendations}, this report provides a prioritized list of behavior modifications ranked by effectiveness:
\begin{enumerate}[leftmargin=*, nosep]
    \item Reducing CPU background processes yields the largest improvement (+42.5\% battery life for 30\% CPU reduction)
    \item Reducing screen brightness from 100\% to 50\% extends battery life by approximately 18\%
    \item Enabling dark mode on OLED screens adds +28\% in specific applications
    \item Maintaining SOC between 20\%--80\% reduces annual capacity loss by approximately 50\%
\end{enumerate}

\textbf{For OS designers}, we propose adaptive power management algorithms including context-aware CPU scheduling, intelligent display dimming, and proactive notification batching.

\textbf{Eventually}, robustness and sensitivity analysis of the model are tested. The model demonstrates stability under parameter perturbations: $\pm$10\% variation in nominal capacity results in only $\pm$6\% change in predicted battery life. When validated against real smartphone discharge data, the model achieves time-to-empty prediction accuracy within $\pm$30 minutes (RMSE = 0.35 hours) for typical usage scenarios ($P$ = 1.5--2.5 W, $T$ = 20--30\textdegree C).

\vspace{0.5cm}
\textbf{Keywords:} Smartphone Battery, State of Charge (SOC), Equivalent Circuit Model (ECM), Ordinary Differential Equation, Time-to-Empty Prediction, Battery Aging, Sobol Sensitivity Analysis, Monte Carlo Simulation
